% Lab 4.1 report.
%
% HOW TO USE
%   1. Save your screenshots into ../screenshots/ using the file names shown
%      in the red placeholder boxes (.png, .jpg or .jpeg all work).
%   2. Fill in everything printed in red (search this file for \fillin).
%   3. Build with `make report` from the project root -> report/report.pdf
%
% A screenshot that is not there yet shows up as a red placeholder box, so
% the document always compiles and you can see what is still missing.
\documentclass[11pt,a4paper]{article}

\usepackage[T1]{fontenc}
\usepackage[utf8]{inputenc}
\usepackage{lmodern}
\usepackage[margin=2.5cm]{geometry}
\usepackage{graphicx}
\usepackage[export]{adjustbox}
\usepackage{xcolor}
\usepackage{booktabs}
\usepackage{float}
\usepackage{listings}
\usepackage[hidelinks]{hyperref}

% ---------------------------------------------------------------- details
\newcommand{\studentname}{Santiago Soto}
\newcommand{\studentid}{M00121124}
\newcommand{\course}{CS 375 Operating Systems}

% ----------------------------------------------------------------- macros
\newcommand{\shotdir}{../screenshots/}

\newcommand{\fillin}[1]{\textcolor{black}{[\,#1\,]}}

% Never enlarged beyond its natural size, so small captures stay sharp.
\newcommand{\shotimage}[1]{%
  \includegraphics[max width=\linewidth,max height=0.40\textheight]{#1}}

\newcommand{\shotfigure}[2]{%
  \begin{figure}[H]\centering #1\caption{#2}\end{figure}}

% \shot{file-name-without-extension}{caption}
% Includes the screenshot if it exists, otherwise a placeholder box.
\newcommand{\shot}[2]{%
  \IfFileExists{\shotdir#1.png}{\shotfigure{\shotimage{\shotdir#1.png}}{#2}}{%
  \IfFileExists{\shotdir#1.jpg}{\shotfigure{\shotimage{\shotdir#1.jpg}}{#2}}{%
  \IfFileExists{\shotdir#1.jpeg}{\shotfigure{\shotimage{\shotdir#1.jpeg}}{#2}}{%
  \typeout{MISSING-SHOT #1}\shotfigure{\fbox{\parbox[c][3.2cm][c]{0.9\linewidth}{\centering\color{red}%
    Missing screenshot\\[4pt]\texttt{screenshots/\detokenize{#1}.png}}}}{#2}}}}}

% \shotopt{file}{caption}: same, but silently skipped when the file is absent.
\newcommand{\shotopt}[2]{%
  \IfFileExists{\shotdir#1.png}{\shotfigure{\shotimage{\shotdir#1.png}}{#2}}{%
  \IfFileExists{\shotdir#1.jpg}{\shotfigure{\shotimage{\shotdir#1.jpg}}{#2}}{%
  \IfFileExists{\shotdir#1.jpeg}{\shotfigure{\shotimage{\shotdir#1.jpeg}}{#2}}{}}}}

\newcommand{\code}[1]{\texttt{\detokenize{#1}}}

\lstset{
  language=C++,
  basicstyle=\ttfamily\small,
  keywordstyle=\color{blue!60!black},
  commentstyle=\color{green!40!black},
  stringstyle=\color{red!60!black},
  frame=single,
  breaklines=true,
  columns=fullflexible,
  showstringspaces=false,
}

\title{Lab 4.1\\Client-Server Applications in C++:\\From Echo Server to Simple Web Server}
\author{\studentname\\\studentid\\\course}
\date{\today}

\begin{document}
\maketitle
\tableofcontents
\newpage

% =========================================================================
\section{Overview}

This lab builds a network server in four steps. Each step keeps the previous
one's socket code and changes only how clients are scheduled:

\begin{center}
\begin{tabular}{@{}cll@{}}
\toprule
Part & Program & Concurrency model \\
\midrule
1 & \code{echo_server} & One client at a time \\
2 & \code{multi_threaded_server} & One thread per client \\
3 & \code{thread_pool_server} & Fixed pool of 10 workers and a task queue \\
4 & \code{http_server} & The Part 3 pool, serving HTTP instead of echo \\
\bottomrule
\end{tabular}
\end{center}

Code shared by every program (socket setup, logging, argument validation and
error handling) lives in \code{src/common.hpp}, so each part's source file
contains only what is new in that part.

\paragraph{Environment.}
Ubuntu 22.04.3 LTS on WSL2, g++ 11.4.0, 16 CPU cores. Built with
\code{g++ -std=c++17 -Wall -Wextra -O2 -pthread}.

\paragraph{Requirements common to all parts.}
\begin{itemize}
  \item \textbf{Error handling.} Ports are validated (whole number, 1--65535).
        Failures of \code{socket}, \code{bind}, \code{listen}, \code{accept}
        and \code{connect} are reported with the reason; a port that is
        already in use and a refused connection each get a specific hint.
        Partial writes, interrupted system calls and clients that disconnect
        early are handled without crashing the server.
  \item \textbf{Logging.} Every event is written to the console with a
        millisecond timestamp. The logger holds a mutex, so lines from
        different threads never mix inside a line.
  \item \textbf{Simulated work.} The echo servers take an optional
        \code{delay_ms} argument. Echoing a few bytes takes microseconds,
        which makes queuing and concurrency invisible; the delay stands in
        for real work so the difference between the designs can be measured.
\end{itemize}

% =========================================================================
\section{Part 1: Simple Echo Server}

The server accepts a client, echoes whatever it sends, closes the connection
and only then accepts the next client.

\subsection{Ten sample requests}

The script \code{scripts/test_echo.sh} sends ten messages and compares each
echo with what was sent.

\begin{center}
\begin{tabular}{@{}rll@{}}
\toprule
\# & Request & Result \\
\midrule
1 & Plain message & Echoed correctly (25 bytes) \\
2 & Longer sentence & Echoed correctly (74 bytes) \\
3 & Empty string & Handled; 0 bytes sent and received \\
4 & 3000 bytes (larger than the 1024-byte buffer) & Echoed correctly in 3 chunks \\
5 & Multi-line message & Echoed correctly (39 bytes) \\
6 & Special characters & Echoed correctly (46 bytes) \\
7 & Unicode (multi-byte characters) & Echoed correctly (48 bytes) \\
8 & Digits & Echoed correctly (21 bytes) \\
9 & Exactly 1024 bytes & Echoed correctly in 1 chunk \\
10 & Final message & Echoed correctly (20 bytes) \\
\bottomrule
\end{tabular}
\end{center}

\shot{p1_requests}{Part 1, server (\texttt{./bin/echo\_server 8080}): log of the ten sample requests.}
\shotopt{p1_requests_2}{Part 1, client (\texttt{scripts/test\_echo.sh 8080}): each echo matches what was sent.}
\shot{p1_telnet}{Part 1, server: the line received from the telnet session.}
\shotopt{p1_telnet_2}{Part 1, client (\texttt{telnet localhost 8080}): the typed line is echoed back.}

\subsection{Extension: multi-line and long messages}

The starting code called \code{read} once with a 1024-byte buffer. TCP is a
byte stream, not a sequence of messages: one \code{read} returns whatever has
arrived so far, so anything beyond the first 1024 bytes was never echoed. The
server now reads in a loop and echoes each chunk until \code{read} returns 0
(end of file). The server log for request 4 shows the 3000-byte message
arriving as several chunks.

The client marks the end of its message by half-closing the connection with
\code{shutdown(fd, SHUT_WR)}. That delivers end-of-file to the server while
leaving the other direction open for the echo.

The starting client also converted the host with \code{inet_pton}, which
accepts only numeric addresses, so \code{./client localhost 8080 ...} could
not work as written. The client uses \code{getaddrinfo} instead, which
resolves names as well as addresses.

\subsection{Queuing}

With a 2000\,ms delay per client, three clients were started at the same
moment.

\shot{p1_queuing}{Part 1, server (\texttt{./bin/echo\_server 8080 2000}): three simultaneous clients are served one after another, two seconds apart.}
\shotopt{p1_queuing_2}{Part 1, client (\texttt{scripts/run\_clients.sh 8080 3}): total time for the three clients.}

All three connections are established immediately, because the kernel
completes the TCP handshake and holds finished connections in the listen
backlog. The server, however, is busy with the first client and does not call
\code{accept} again until it is done. The timestamps show the clients being
served about two seconds apart, for a total of 6.0\,s.
One slow client therefore delays everyone behind it.

\subsection{Error handling}

\shot{p1_errors}{Server started with an invalid port (\texttt{abc}) and an out-of-range port (\texttt{99999}).}
\shotopt{p1_errors_2}{Client run while no server is listening: connection refused.}

% =========================================================================
\section{Part 2: Multi-Threaded Echo Server}

The main thread now only accepts connections. Each client is handed to a new
detached \code{pthread}, so a slow client no longer blocks the others.

\subsection{Concurrent requests}

\shot{p2_concurrent}{Part 2, server: ten parallel clients. The log lines of different clients interleave.}
\shotopt{p2_concurrent_2}{Part 2, client: output of the ten parallel clients.}

In Part 1 every client's ``connected / received / disconnected'' lines
appeared as one block. Here all ten clients connect within a few milliseconds
and their lines are interleaved, which shows that they are being served at
the same time.

\subsection{Performance compared with Part 1}

Ten clients, 500\,ms of simulated work per client:

\begin{center}
\begin{tabular}{@{}lrr@{}}
\toprule
Server & Clients & Total time \\
\midrule
Part 1, single-threaded & 10 & 5.02\,s \\
Part 2, thread per client & 10 & 0.55\,s \\
\bottomrule
\end{tabular}
\end{center}

\shot{p2_timing_part1}{Part 1 under the timing test.}
\shot{p2_timing_part2}{Part 2 under the same test.}

Part 1 takes roughly $10 \times 0.5$\,s because the delays add up. In Part 2
the ten delays overlap, so the total is close to a single delay: about nine
times faster for the same work.

\subsection{Shared connection counter}

The total number of connections and the number of active threads are shared
by all threads and protected by a \code{pthread_mutex_t}. An increment is a
read, an add and a write; without the mutex two threads can read the same old
value and one update is lost. The lock makes the three steps indivisible.

\subsection{Discussion: thread creation overhead}

Each connection pays for creating and destroying a thread, and each live
thread reserves its own stack (8\,MiB of address space by default on Linux)
and adds scheduling work. For ten clients this is negligible. The danger is
that the number of threads is decided by the clients, not the server: several
thousand simultaneous connections would exhaust memory or the thread limit,
\code{pthread_create} would fail, and a client can cause this deliberately as
a denial-of-service attack.

Two measures are implemented. The return value of \code{pthread_create} is
checked and the connection is closed cleanly if it fails, and the server
refuses new clients once 200 threads are active (\code{MAX_THREADS}) rather
than growing without bound. Part 3 removes the problem at its source.

% =========================================================================
\section{Part 3: Thread-Pool Echo Server}

Ten worker threads are created once at start-up. The main thread accepts a
connection and pushes it onto a \code{std::queue}; a worker pops it and
serves it. The number of threads no longer depends on the number of clients.

\subsection{Queuing under load}

Twenty clients were launched against the ten workers, with message sizes
varying from 20 bytes to several kilobytes.

\shot{p3_queue_sizes}{Part 3, server, queue lines only: clients 11--20 are queued while all ten workers are busy, then dequeued as workers become free.}
\shotopt{p3_queue}{Part 3, server, full log: the last clients leave the queue and the second batch is served by the same ten workers about 500\,ms after the first.}
\shotopt{p3_queue_2}{Part 3, client (\texttt{scripts/run\_clients.sh 8080 20}): all twenty echoes match.}
\shotopt{p3_queue_time}{Part 3, client: total time for the twenty clients.}

The first ten clients are picked up immediately. The next ten wait and the
logged queue size climbs to 10. As
workers finish they dequeue the waiting clients and the size falls back to
zero. Twenty clients took 1.03\,s: two rounds of ten at 500\,ms each.

\subsection{Synchronisation}

The mutex protects the queue, which the main thread and all workers modify.
The condition variable lets idle workers sleep until there is work: the main
thread calls \code{notify_one} after each push, waking exactly one worker.
Without it the workers would have to poll the queue in a loop and burn CPU
while idle. The wait uses a predicate (queue not empty, or pool stopping), so
spurious wake-ups are harmless.

\subsection{Graceful shutdown}

\shot{p3_shutdown}{Part 3, server: Ctrl+C while clients are still being served. Work already accepted is finished before the server exits.}
\shotopt{p3_shutdown_2}{Part 3, server: shutdown, continued.}

A \code{SIGINT}/\code{SIGTERM} handler only sets a flag. The handler is
installed without \code{SA_RESTART}, so the blocked \code{accept} returns
\code{EINTR} and the main loop sees the flag. The signals are blocked while
the workers are created, which guarantees delivery to the main thread. The
server then stops accepting, sets \code{stop_pool}, wakes every worker and
joins them. A worker exits only when the pool is stopping \emph{and} the
queue is empty, so every client that was already accepted is still served;
the screenshot shows the queue being drained before ``Shutdown complete''.

\subsection{Resource use compared with Part 2}

Both servers were given twenty clients with a 15\,000\,ms delay and inspected
with \code{ps} while the clients were connected.

\begin{center}
\begin{tabular}{@{}lrrrr@{}}
\toprule
Server & Clients & Threads (NLWP) & Memory (RSS) & CPU \\
\midrule
Part 2, thread per client & 20 & 21 & 4348\,kB & 0.0\,\% \\
Part 3, pool of 10 & 20 & 11 & 4068\,kB & 0.0\,\% \\
\bottomrule
\end{tabular}
\end{center}

\shot{p3_top_threads}{Part 2 under load: 21 threads, one per connected client plus the main thread.}
\shot{p3_top_pool}{Part 3 under load: 11 threads, the same as when idle.}

Part 2's thread count follows the number of clients, while Part 3 stays at
eleven (ten workers and the main thread) regardless of load. CPU use is zero
in both because the threads are asleep in the simulated delay, and resident
memory differs by only 280\,kB: a thread's stack is reserved address space
that costs real memory only once it is used, and these threads do very little.
At twenty clients the difference is therefore small. What matters is how it
grows: Part 2 adds a thread for every further client without limit, and Part 3
adds none. The price is
latency: when all workers are busy, a new client waits in the queue. That is
the better failure mode, because an overloaded pool becomes slower while an
overloaded thread-per-client server runs out of resources.

% =========================================================================
\section{Part 4: Simple HTTP Web Server}

The Part 3 pool is reused unchanged; the workers now parse an HTTP request
and send a file from \code{./www}.

\subsection{Request parsing and response}

\code{parse_request} reads until the blank line that ends the headers (a
request may arrive in several pieces), then:
\begin{itemize}
  \item checks that the request line is exactly \emph{method, target,
        version}, with an \code{HTTP/1.x} version and a target starting with
        \code{/};
  \item parses the headers into a map with lower-cased names;
  \item splits the target into path and query string, percent-decodes the
        path (\code{\%20} becomes a space) and parses the query parameters;
  \item reads a POST body of exactly \code{Content-Length} bytes.
\end{itemize}
\code{GET} and \code{HEAD} serve files, \code{POST} echoes the body it
received, and any other method is answered with \code{405}. Malformed
requests receive \code{400}, oversized headers \code{431}, oversized bodies
\code{413} and clients that stall \code{408}. A five-second socket timeout
stops an idle connection from occupying a worker indefinitely. The full
request (request line, headers, query parameters and body) is written to the
log.

Files are read in binary mode so that images are not altered, and the
response carries the exact \code{Content-Length}. The \code{Content-Type} is
chosen from the file extension. This matters in practice: a browser will not
apply a stylesheet that is not served as \code{text/css}, and needs
\code{image/png} to render the logo.

\subsection{Ten sample requests}

\begin{center}
\begin{tabular}{@{}rllll@{}}
\toprule
\# & Request & Expected & Status & Content-Type \\
\midrule
1 & \code{/} & 200 & 200 & text/html \\
2 & \code{/index.html} & 200 & 200 & text/html \\
3 & \code{/style.css} & 200 & 200 & text/css \\
4 & \code{/logo.png} & 200 & 200 & image/png \\
5 & \code{/subpage.html} & 200 & 200 & text/html \\
6 & \code{/missing.txt} & 404 & 404 & text/html (error page) \\
7 & \code{/favicon.ico} & 404 & 404 & text/html (error page) \\
8 & \code{/..} & 404 & 404 & text/html (error page) \\
9 & \code{/?test=1} & 200 & 200 & text/html \\
10 & \code{/very/long/path/to/file.html} & 404 & 404 & text/html (error page) \\
\bottomrule
\end{tabular}
\end{center}

\shot{p4_tests}{Part 4: all sample requests and traversal attacks, with status, type and size.}

\subsection{Browser}

\shot{p4_browser_index}{The index page with its stylesheet and logo applied.}
\shot{p4_browser_subpage}{The subpage.}
\shot{p4_browser_404}{The 404 page for the missing link.}
\shotopt{p4_browser_log}{Server log while the browser loads the index page: full requests are logged and handled by several workers.}

Loading the index page makes the browser send several requests at once (the
page, the stylesheet, the logo and a favicon), which the server log shows
being handled by different workers.

\subsection{curl}

\shot{p4_curl_index}{\texttt{curl -v /}: request and response headers, then the HTML.}
\shot{p4_curl_css}{\texttt{curl -v /style.css}: served as \texttt{text/css}.}
\shot{p4_curl_png}{\texttt{curl -v /logo.png}: served as \texttt{image/png}.}
\shot{p4_curl_404}{\texttt{curl -v /missing.txt}: 404 Not Found.}
\shotopt{p4_curl_indexhtml}{\texttt{curl -v /index.html}.}
\shotopt{p4_curl_subpage}{\texttt{curl -v /subpage.html}.}
\shotopt{p4_curl_favicon}{\texttt{curl -v /favicon.ico}: 404 Not Found.}
\shotopt{p4_curl_dotdot}{\texttt{curl -v --path-as-is /..}: 404 Not Found.}
\shotopt{p4_curl_query}{\texttt{curl -v /?test=1}: the query string is parsed and the index page is served.}
\shotopt{p4_curl_longpath}{\texttt{curl -v /very/long/path/to/file.html}: 404 Not Found.}

% =========================================================================
\section{Cybersecurity Challenge: Path Traversal}

\subsection{The attack}

The server builds a file name by appending the URL path to the web root. A
path containing \code{../} walks back out of that directory:
\code{./www/../../../etc/passwd} is \code{/etc/passwd}. The starting code
rejected any path containing ``\code{..}'', but it ran that check on the raw
URL. In \code{/../\%2e\%2e/\%2e\%2e/etc/passwd} the dots are written as
\code{\%2e}, so a check for the two characters ``\code{..}'' finds nothing in
the encoded segments, and they turn into \code{..} only when the path is
decoded later. More generally, searching for known-bad strings is fragile:
it has to anticipate every encoding, and it cannot detect a symbolic link
inside the web root that points outside it, which involves no ``\code{..}''
at all.

\subsection{The defence}

Instead of looking for bad input, the server works out where the request
really leads and checks the result:

\begin{enumerate}
  \item \textbf{Decode once, first.} The path is percent-decoded before any
        check, so every later step sees what the file system will see.
        Malformed escapes and \code{\%00} are rejected with \code{400}.
  \item \textbf{Canonicalise.} \code{std::filesystem::weakly_canonical} turns
        the path into an absolute one with \code{.}, \code{..} and symbolic
        links resolved.
  \item \textbf{Verify containment.} The resolved path must lie inside the
        canonical web root, compared component by component. A plain string
        prefix test would be wrong here: it would accept the directory
        \code{/srv/www-secret} as being inside \code{/srv/www}.
\end{enumerate}

\begin{lstlisting}
fs::path resolved = fs::weakly_canonical(
    g_web_root / fs::path(url_path).relative_path(), ec);
if (ec) return std::nullopt;

auto mismatch = std::mismatch(g_web_root.begin(), g_web_root.end(),
                              resolved.begin(), resolved.end());
if (mismatch.first != g_web_root.end()) {
    lab::log("SECURITY: blocked path traversal attempt: ...");
    return std::nullopt;   // answered with 404
}
\end{lstlisting}

\subsection{Test}

\shot{p4_traversal}{Traversal attempt with plain \texttt{../} segments: 404 Not Found.}
\shotopt{p4_traversal_2}{Traversal attempt with percent-encoded dots (\texttt{\%2e\%2e}): 404 Not Found.}
\shot{p4_traversal_log}{The server log records each blocked attempt and where it would have led.}

curl needs \code{--path-as-is} for this test. By default it removes
\code{../} segments from the URL before sending, so without the flag the
attack never reaches the server and the test proves nothing.

The server answers \code{404} rather than \code{403}, so the response does
not reveal whether the target exists. The script \code{scripts/test_http.sh}
additionally tests encoded slashes (\code{..\%2f}), a path that starts inside
the root and then climbs out, a NUL byte and a malformed escape.
All seven cases passed: the five traversal attempts returned \code{404} and
the two malformed URLs returned \code{400}.

\subsection{Why this matters}

A traversal flaw lets anyone who can reach the server read any file the
server process can read: \code{/etc/passwd}, application source code,
configuration files, database passwords, API keys and private keys. Leaked
credentials usually turn an information leak into a full compromise. The
attack needs no authentication and nothing more than a URL, so it is easy to
automate and is scanned for constantly.

Path traversal is CWE-22 and belongs to \textbf{A01: Broken Access Control},
the first category of the OWASP Top 10 (2021): the server fails to enforce
the rule that clients may only read files inside the web root. The general
lessons are to validate the outcome rather than the input text, to normalise
data before checking it, and to run the server with as few file permissions
as possible so that a mistake exposes less.

% =========================================================================
\section{Conclusion}

\begin{center}
\begin{tabular}{@{}lll@{}}
\toprule
Design & Strength & Weakness \\
\midrule
Single-threaded & Simplest; no shared state & One slow client blocks all others \\
Thread per client & Clients fully concurrent & Thread count set by the clients \\
Thread pool & Bounded resources & Clients wait when the pool is busy \\
\bottomrule
\end{tabular}
\end{center}

The measurements match the designs. With 500\,ms of work per client, the
single-threaded server needed 5.02\,s for ten clients because every client
waited for all the ones ahead of it. One thread per client brought that down
to 0.55\,s, and the pool of ten served twenty clients in 1.03\,s while using
11 threads where the thread-per-client server needed 21.

For a real server I would choose the thread pool. It gives nearly all of the
concurrency of one thread per client, but its resource use is decided by the
server rather than by whoever connects to it, and under overload it slows down
instead of failing. The pool size and the queue then become the two settings
to tune for the workload.

The security exercise showed that the networking code was the easy part.
The original traversal check looked reasonable and was still bypassed by
encoding two characters. Checking where a request actually ends up, rather
than what its text looks like, is the approach I would reuse elsewhere.

\end{document}
